\documentclass[11pt,a4paper]{jlreq}
\usepackage{amsmath,amssymb}
\usepackage[dvipdfmx]{graphicx}
\def\pgfsysdriver{pgfsys-dvipdfmx.def}
\usepackage{tikz}
\usepackage{circuitikz}
\usetikzlibrary{circuits.logic.US,calc}
\usepackage{array}
\usepackage{float}

\begin{document}

\begin{center}
{\LARGE \textbf{情報科学基礎A 第5回 レポート}}
\end{center}
\vspace{1em}

\begin{flushright}
提出日：\today \\
学籍番号：2611193 \\
氏名：中井平蔵
\end{flushright}

\section*{Q1}
\subsection*{(1) 真理値表}
真理値表は以下のようになる。
\vspace{0.4\baselineskip}
\begin{table}[H]
\centering
\caption{3入力多数決回路の真理値表}
\renewcommand{\arraystretch}{1.25}
\begin{tabular}{|>{\centering\arraybackslash}m{1.8cm}|>{\centering\arraybackslash}m{1.8cm}|>{\centering\arraybackslash}m{1.8cm}|>{\centering\arraybackslash}m{1.8cm}|}
\hline
a & b & c & y \\
\hline
0 & 0 & 0 & 0 \\
 \hline
0 & 0 & 1 & 0 \\
 \hline
0 & 1 & 0 & 0 \\
 \hline
0 & 1 & 1 & 1 \\
 \hline
1 & 0 & 0 & 0 \\
 \hline
1 & 0 & 1 & 1 \\
 \hline
1 & 1 & 0 & 1 \\
 \hline
1 & 1 & 1 & 1 \\
\hline
\end{tabular}
\renewcommand{\arraystretch}{1.0}
\end{table}
\vspace{1.0em}

\subsection*{(2) 論理式の導出}
$y=1$ となるのは、$(a,b,c)=(0,1,1),(1,0,1),(1,1,0),(1,1,1)$ のときである。
標準形（加法標準形, Sum of Products）は、真理値表で $y=1$ となる各行に対して
「0の変数は否定、1の変数はそのまま」掛け合わせた最小項を作り、それらをすべてORで足し合わせて作る。
よって、
\[
\begin{aligned}
(0,1,1)&\rightarrow \bar{a}bc,\\
(1,0,1)&\rightarrow a\bar{b}c,\\
(1,1,0)&\rightarrow ab\bar{c},\\
(1,1,1)&\rightarrow abc
\end{aligned}
\]
より標準形は以下のようになる。
\[
y=\bar{a}bc+a\bar{b}c+ab\bar{c}+abc
\]
ここで、$x + x = x$より、$abc = abc + abc + abc$が成り立つ。よって、
\[
\begin{aligned}
y
&=(ab\bar{c}+abc)+(a\bar{b}c+abc)+(\bar{a}bc+abc)\\
&=ab(\bar{c}+c)+ac(\bar{b}+b)+bc(\bar{a}+a)\\
\end{aligned}
\]
ここで、$x + \bar{x} = 1$ より、$\bar{c}+c=1$、$\bar{b}+b=1$、$\bar{a}+a=1$ が成り立つ。\\
したがって、論理式は以下のようになる。
\[
\fbox{$y=ab+bc+ca$}
\]

\section*{Q2}
\noindent
入力ボタン $a,b$、リセット $r$、クロック $c$ を持つシンプルなデジタル錠を設計する。
解錠条件は「Aを押した次のクロックでBを押す」のみである。

\subsection*{(1) 状態定義}
\begin{itemize}
\item $S_0$: Initial（初期ロック状態）
\item $S_1$: A\_pressed（Aが押されたことを記憶した状態）
\item $S_2$: Unlocked（解錠状態）
\end{itemize}

\subsection*{(2) 状態割当表と遷移表}
状態を次の表のように2ビットで表す。
\vspace{0.4\baselineskip}
\begin{table}[H]
\centering
\caption{状態割当表}
\renewcommand{\arraystretch}{1.25}
\begin{tabular}{|>{\centering\arraybackslash}m{4.0cm}|>{\centering\arraybackslash}m{4.0cm}|}
\hline
State & $(Q_1^{(n)}Q_0^{(n)})$ \\
\hline
$S_0$ & 00 \\
\hline
$S_1$ & 01 \\
\hline
$S_2$ & 10 \\
\hline
\end{tabular}
\renewcommand{\arraystretch}{1.0}
\end{table}

\vspace{0.5\baselineskip}

入力と状態遷移表の関係は以下のようになる。

\vspace{0.6\baselineskip}

\begin{table}[H]
\centering
\caption{状態遷移表}
\renewcommand{\arraystretch}{1.2}
\begin{tabular}{|>{\centering\arraybackslash}m{1.2cm}|>{\centering\arraybackslash}m{1.2cm}|>{\centering\arraybackslash}m{1.0cm}|>{\centering\arraybackslash}m{1.0cm}|>{\centering\arraybackslash}m{1.0cm}|>{\centering\arraybackslash}m{1.2cm}|>{\centering\arraybackslash}m{1.2cm}|>{\centering\arraybackslash}m{1.0cm}|}
\hline
$Q_1^{(n)}$ & $Q_0^{(n)}$ & $a$ & $b$ & $r$ & $Q_1^{(n+1)}$ & $Q_0^{(n+1)}$ & Out \\
\hline
0 & 0 & 0 & 0 & 0 & 0 & 0 & 0 \\
\hline
0 & 0 & 1 & 0 & 0 & 0 & 1 & 0 \\
\hline
0 & 0 & 0 & 1 & 0 & 0 & 0 & 0 \\
\hline
0 & 0 & 0 & 0 & 1 & 0 & 0 & 0 \\
\hline
0 & 1 & 0 & 0 & 0 & 0 & 0 & 0 \\
\hline
0 & 1 & 1 & 0 & 0 & 0 & 1 & 0 \\
\hline
0 & 1 & 0 & 1 & 0 & 1 & 0 & 0 \\
\hline
0 & 1 & 0 & 0 & 1 & 0 & 0 & 0 \\
\hline
1 & 0 & 0 & 0 & 0 & 1 & 0 & 1 \\
\hline
1 & 0 & 1 & 0 & 0 & 1 & 0 & 1 \\
\hline
1 & 0 & 0 & 1 & 0 & 1 & 0 & 1 \\
\hline
1 & 0 & 0 & 0 & 1 & 0 & 0 & 0 \\
\hline
\end{tabular}
\renewcommand{\arraystretch}{1.0}
\end{table}

\vspace{0.6\baselineskip}

\subsection*{(3) 論理式}
\begin{itemize}
\item $Q_1^{(n+1)}$ の導出 \\
$Q_1^{(n+1)}=1$ となるのは、\\
\[
(Q_1^{(n)},Q_0^{(n)},a,b,r)=(0,1,0,1,0),(1,0,0,0,0),(1,0,1,0,0),(1,0,0,1,0)
\]
よって、
\[
Q_1^{(n+1)}=\overline{Q_1^{(n)}}Q_0^{(n)}\overline{a}b\overline{r}+Q_1^{(n)}\overline{Q_0^{(n)}}\overline{a}\,\overline{b}\,\overline{r}+Q_1^{(n)}\overline{Q_0^{(n)}}a\overline{b}\,\overline{r}+Q_1^{(n)}\overline{Q_0^{(n)}}\overline{a}b\overline{r}
\]
初めに、2〜4項を $\overline{r}\,Q_1^{(n)}\overline{Q_0^{(n)}}$ でくくると、\\
\[
Q_1^{(n+1)}=\overline{r}\,\overline{Q_1^{(n)}}Q_0^{(n)}\overline{a}b+\overline{r}\,Q_1^{(n)}\overline{Q_0^{(n)}}(\overline{a}\,\overline{b}+a\overline{b}+\overline{a}b)
\]
ここで、$\overline{a}\,\overline{b}+a\overline{b}=(\overline{a}+a)\overline{b}=\overline{b}$より、
\vspace{0.3\baselineskip}
\[
Q_1^{(n+1)}=\overline{r}\,\overline{Q_1^{(n)}}Q_0^{(n)}\overline{a}b+\overline{r}\,Q_1^{(n)}\overline{Q_0^{(n)}}(\overline{b}+\overline{a}b)
\]
また、$\overline{b}+\overline{a}b$ は場合分けで次のように分かる。
\begin{itemize}
\item $b=0$ のとき：$\overline{b}+\overline{a}b=1+0=1$
\item $b=1$ のとき：$\overline{b}+\overline{a}b=0+\overline{a}=\overline{a}$
\end{itemize}
したがって、$\overline{b}+\overline{a}b=\overline{a}+\overline{b}$ である。\\
よって、
\vspace{0.3\baselineskip}
\[
Q_1^{(n+1)}=\overline{r}\,\overline{Q_1^{(n)}}Q_0^{(n)}\overline{a}b+\overline{r}\,Q_1^{(n)}\overline{Q_0^{(n)}}(\overline{a}+\overline{b})
\]
さらに、$a,b$ は同時に $1$ にならないため、$\overline{a}+\overline{b}=1$である。よって、 \\
\[
Q_1^{(n+1)}=\overline{r}\,\overline{Q_1^{(n)}}Q_0^{(n)}\overline{a}b+\overline{r}\,Q_1^{(n)}\overline{Q_0^{(n)}}
\]
また、$(Q_1^{(n)},Q_0^{(n)}) = (1,1)$となることは無いため、$Q_1^{(n)}\overline{Q_0^{(n)}}=Q_1^{(n)}$ である。よって、\\
\[
Q_1^{(n+1)}=\overline{r}\,\overline{Q_1^{(n)}}Q_0^{(n)}\overline{a}b +\overline{r}\,Q_1^{(n)}
\]
同様に、$(Q_1^{(n)},Q_0^{(n)}) = (1,1)$となることは無いため、$\overline{Q_1^{(n)}}Q_0^{(n)}=Q_0^{(n)}$ である。よって、\\
\[
Q_1^{(n+1)}=\overline{r}\,Q_0^{(n)}\overline{a}b +\overline{r}\,Q_1^{(n)}
\]
さらに、$a,b$ は同時に $1$ にならないため、$\overline{a}b = b$である。よって、 \\
\[
Q_1^{(n+1)}=\overline{r}\,\bigl(Q_1^{(n)}+Q_0^{(n)}b\bigr)
\]
したがって、
\vspace{0.3\baselineskip}
\[
Q_1^{(n+1)}=\overline{r}\,\bigl(Q_1^{(n)}+Q_0^{(n)}b\bigr)
\]

\item $Q_0^{(n+1)}$ の導出 \\
$Q_0^{(n+1)}=1$ となるのは、\\
\[
(Q_1^{(n)},Q_0^{(n)},a,b,r)=(0,0,1,0,0),(0,1,1,0,0)
\]
よって、
\[
Q_0^{(n+1)}=\overline{Q_1^{(n)}}\,\overline{Q_0^{(n)}}a\overline{b}\,\overline{r}+\overline{Q_1^{(n)}}Q_0^{(n)}a\overline{b}\,\overline{r}
\]
初めに、$\overline{r}\,\overline{Q_1^{(n)}}a\overline{b}$ でくくると、\\
\[
Q_0^{(n+1)}=\overline{r}\,\overline{Q_1^{(n)}}a\overline{b}\,(\overline{Q_0^{(n)}}+Q_0^{(n)})
\]
\vspace{-0.4\baselineskip}
$\overline{Q_0^{(n)}}+Q_0^{(n)}=1$ より、\\
\[
Q_0^{(n+1)}=\overline{r}\,\overline{Q_1^{(n)}}a\overline{b}
\]
\end{itemize}

\begin{itemize}
\item 出力 \\
出力が1となるのは、\\
\[
(Q_1^{(n)},Q_0^{(n)},a,b,r)=(1,0,0,0,0),(1,0,1,0,0),(1,0,0,1,0)
\]
よって、
\[
\mathrm{Out}
=Q_1^{(n)}\overline{Q_0^{(n)}}\overline{a}\,\overline{b}\,\overline{r}
+Q_1^{(n)}\overline{Q_0^{(n)}}a\overline{b}\,\overline{r}
+Q_1^{(n)}\overline{Q_0^{(n)}}\overline{a}b\overline{r}
\]
初めに、$Q_1^{(n)}\overline{Q_0^{(n)}}\overline{r}$ でくくると、\\
\[
\mathrm{Out}=Q_1^{(n)}\overline{Q_0^{(n)}}\overline{r}\,(\overline{a}\,\overline{b}+a\overline{b}+\overline{a}b)
\]
$\overline{a}\,\overline{b}+a\overline{b}=\overline{b}$、さらに
$\overline{b}+\overline{a}b=\overline{a}+\overline{b}$ より、
\[
\mathrm{Out}=Q_1^{(n)}\overline{Q_0^{(n)}}\overline{r}\,(\overline{a}+\overline{b})
\]
ここで、$a,b$ は同時に1にならないため $\overline{a}+\overline{b}=1$ である。したがって、
\[
\mathrm{Out}=Q_1^{(n)}\overline{Q_0^{(n)}}\overline{r}
\]
\end{itemize}
以上より、
\[
Q_1^{(n+1)}=\overline{r}\,\bigl(Q_1^{(n)}+Q_0^{(n)}b\bigr),\quad
Q_0^{(n+1)}=\overline{r}\,\overline{Q_1^{(n)}}a\overline{b},\quad
\mathrm{Out}=Q_1^{(n)}\overline{Q_0^{(n)}}\overline{r}
\]

\subsection*{(4) 回路構成（D-FF）}
次状態論理を以下のようにおき、2個のD-FFで状態 $(Q_1,Q_0)$ を保持する。
\[
D_1=Q_1^{(n+1)},\quad D_0=Q_0^{(n+1)}
\]

\begin{figure}[H]
\centering
\noindent
% 画像配置調整用パラメータ
% width: 表示サイズ, trim: left bottom right top の順（単位は pt）
\includegraphics[
  width=0.44\textheight,
  angle=90,
  origin=c,
  trim=450 450 450 450,
  clip
]{img/IMG_9454.pdf}
\vspace{-2.4\baselineskip}
\setlength{\abovecaptionskip}{2pt}
\caption{回路構成図（D-FFを用いた実装）}
\end{figure}

\end{document}
